\section{Introduction}
An exact rational periodic locus asks more than whether periodic points
are finite or their periods are bounded. For a polynomial automorphism
over a function field, a local symbolic description can allow many
itineraries that do not come from rational functions. We determine this
rationality restriction for the entire family
\begin{equation}\label{eq:map}
 H_{a,c}(x,y)=(y,y^2+c-a x),\qquad
 a\in k^*,\quad c\in k[t]\setminus k,\quad \ch k\ne2.
\end{equation}
The domain throughout is $k(t)^2$. In particular, the nonconstant
parameter is polynomial, while the determinant is constant.

The reduction is elementary but global. At a finite polynomial prime,
a coordinate with maximal pole order contradicts the quadratic
recurrence. At infinity, all coordinates in all periodic orbits have
the same positive degree. Factoring the difference of their squares
then puts every coordinate in $\{P+b,-P+b:b\in k\}$ for one common
nonconstant polynomial $P$. The recurrence determines the offsets from
neighboring signs. Only sixteen signed labels remain, and their exact
transition graph decides which are recurrent.

The main conclusions are the following.
\begin{enumerate}
\item The parameter must have the unique centered form $c=-P^2+C$,
where uniqueness of $P$ is understood up to sign. A seven-row atlas
then constructs every rational cycle, including all overlaps.
\item The sharp universal total is fourteen rational periodic points.
The unique equality locus is $\ch k=3$, $a=-1$, $C=-1$; the
ordinary periods there are four and five.
\item Outside characteristics two and three, the sharp bound is eight
or six according as $-1$ is or is not a square in $k$. Longer-than-two
periods in any allowed characteristic require $a^4=1$.
\end{enumerate}
No hypothesis of perfection, algebraic closure, finite cardinality,
ordering or local compactness is imposed on $k$.

\paragraph{Arithmetic bounds and conventions.}
Canonical heights and non-Archimedean escape are established tools for
H\'enon arithmetic. Ingram's function-field finiteness and bad-place
bounds are prior results \citep{ingram2014canonical}; neither the
existence of such bounds nor the local escape method is a contribution
here. His normalization is
$\varphi_\alpha(x,y)=(\alpha y,x+f(y))$. For
$L(x,y)=(-a x,y)$ one has
\[
 L H_{a,c}L^{-1}(x,y)=(-a y,x+y^2+c).
\]
Thus the source's special value $\alpha=1$ corresponds to our $a=-1$,
not to the conservative value $a=1$. Our comparison uses the
introduction and relevant local arguments of arXiv:1111.3609v1,
especially its stated quadratic routing; the journal metadata in the
bibliography does not assert that its final proof text was read.
Those accessed statements give finiteness and bounds, not the present
all-determinant equality locus and every signed cycle.

\paragraph{A local horseshoe is not the global rational locus.}
Allen, DeMark and Petsche study
$\varphi_{A,B}(X,Y)=(A+B Y-X^2,X)$ over complete, locally compact
non-Archimedean fields of odd residue characteristic
\citep{allen2018nonarchimedean}. Their horseshoe theorem gives a full
two-symbol bilateral shift in the appropriate region
$|A|>\max(1,|B|^2)$ with $A$ a square.
The change $J(x,y)=(-y,-x)$ gives
$JH_{a,c}J^{-1}=\varphi_{-c,-a}$.
For finite odd $k$ the completion $k((1/t))$ fits their field setting,
but a shift over the completion does not select the itineraries in
$k(t)$. For arbitrary $k$ the completion need not be locally compact.
Our graph is a partial permutation precisely because it imposes this
global rationality restriction. The coding method itself is classical.

\paragraph{Nearby repository results.}
The earlier unpublished repository note C412 on conservative integral
quadratic H\'enon maps already uses sign/constant-offset equations at
$a=1$. We credit that encoding, as well as its real separation and
finite small-parameter completion; C412 is not represented as a journal
publication. The present proof replaces those integer-size arguments
by across-cycle polynomial-degree rigidity and resolves every allowed
constant determinant and characteristic. Constant-coefficient
finite-field height distributions concern a different parameter family
and observable. These distinctions identify the claim studied here;
they are not a claim of an exhaustive priority search.

Section~\ref{sec:atlas} states the atlas and reconstruction. Sections
\ref{sec:rigidity}--\ref{sec:exhaustion} prove the global reduction,
exact graph and all-field exhaustion. Section~\ref{sec:bounds}
retains every row collision to prove the sharp bounds.
